\clearpage
\section{Implementation boundaries and numerical reference}

\subsection{What is exact in the supplied implementation}

The FIFO routine consumes an ordered list. The proportional routine calculates rational entitlements, applies an integer minimum, and allocates residual demand FIFO. Both operate on the ordinary-order state declared in Section~\ref{sec:allocation}. They do not infer TOP status, LMM entitlements, hidden replenishment, or a product's configured sequence. The opening routine covers the five tabulated tie cases and the unique-volume case; no production fallback is supplied for remaining exact ties.

The normal money-value map uses decimal half-away-from-zero rounding at one-contract value level. No special notional rounding, inverse currency conversion, option futures-style premium, or daily-value adjustment is applied. The numerical examples explicitly choose instruments and conventions for which the shown arithmetic is consistent with the declared map.

The margin routine uses exactly the six vectors and three additions in Section~\ref{sec:margin}. The plotting grid is continuous, while the example positions are integer contracts. The collateral example uses one legal account, fractional bond-value allocation, a synthetic haircut, and a cash-only current variation obligation. Extending it across accounts or currencies requires additional constraints.

\begin{table}[htbp]
\centering\small
\begin{tabularx}{\textwidth}{@{}lX@{}}
\toprule
Quantity & Reproduction target\\
\midrule
37-lot allocation & FIFO $(12,25,0)$; PR/FIFO $(5,10,22)$\\
4-lot fragmentation & One event $(2,0,2)$; two 2-lot events $(2,2,0)$\\
ES depth walk & Average 6000.225; shortfall \$112.50\\
Shared implication & Individually $5+6$; jointly at most 8\\
Completion indices & FIFO $(1,2,3)$; PR/FIFO $(3,3,3)$\\
Opening price & 102, with 10 matched contracts\\
VWAP perturbation & 6000.25 to $6000.423076\ldots$; rounded mark 6000.50\\
Three-day variation & \$387.50, $-\$700.00$, \$1,500.00\\
Total cash gain & \$1,187.50; final cash \$13,187.50\\
Illustrative margin & \$11,200 at $(5,-5)$; \$12,800 at $(5,-3)$\\
Peak funding & \$10,000 versus zero; \$15,000 with margin shock\\
Collateral bottleneck & \$23,000 credited value; only \$5,000 current cash\\
Option example & $-\$1,800$ premium; \$3,000 variation; \$1,200 net\\
Delivery comparison & Net costs \$0, $-\$200$, \$1,200; B cheapest at 110\\
Waterfall & Loss 17; draws $(8,2,5,2)$ in synthetic millions\\
\bottomrule
\end{tabularx}
\caption{Principal numerical checkpoints. Each belongs to its stated model; the margin and collateral examples are separate from the fixed-rate integrated ledger.}
\end{table}

\subsection{Source dating and reproducible compilation}

The manuscript date and consultation date are September 11, 2026 (UTC). A 2026 bibliography year for a continuously maintained operational page records the consulted snapshot; it does not imply original publication in that year. The money-calculation specification retains its stated 2015 document year. Source URLs, titles, and consultation notes are supplied in the bibliography and in \texttt{source\_register.md}.

Compile \texttt{The\_Mathematics\_of\_CME.tex} with pdfLaTeX in Overleaf. Included figures require no Python or shell escape. To regenerate them and the numerical ledger, run \texttt{reproduce.py}. The source archive includes a SHA256 manifest.
\FloatBarrier
